\section{Performance Management (KPIs)}
\label{sec:kpi}

% Mandatory: list KPIs AND argue why. ~0.75 page.
We measure four headline KPIs and one secondary KPI, chosen to expose the
constraint and to speak directly to management's goal of maximising output.

\begin{table}[H]
\centering
\caption{Performance metrics used in this study.}
\label{tab:kpis}
\begin{tabularx}{\textwidth}{l l X}
\toprule
\textbf{KPI} & \textbf{Type} & \textbf{Why it was chosen} \\
\midrule
Weekly throughput (good bottles / tablets) & Headline & Directly operationalises the 2030 goal of maximising output. \\
Resource utilisation per station & Headline & Primary bottleneck detector: the station near \SI{100}{\percent} while others starve is the constraint. \\
Batch flow time (release $\rightarrow$ shipped) & Headline & Responsiveness; sensitive to cleaning-team contention and changeovers. \\
WIP / queue length per station & Headline & Localises \emph{where} batches accumulate; corroborates utilisation. \\
Scrap rate / overall yield & Secondary & Links to the Industry~4.0 lever (press scrap \SI{10}{\percent}$\rightarrow$\SI{4}{\percent}). \\
\bottomrule
\end{tabularx}
\end{table}

% Constraint we respect but do not analyse heavily:
The planned-vs-actual completion date (the $\le 1$-day deviation constraint) is
respected as a modelling constraint rather than reported as a headline KPI.
